\documentclass[a4paper,12pt, print, oneside]{book}
\usepackage[T1]{fontenc} %% set font encode
\usepackage[utf8]{inputenc}
\usepackage[unicode,hyperfootnotes=false]{hyperref}
\usepackage{vntex}
\usepackage[vietnamese=nohyphenation]{hyphsubst} % fix warning when using vietnamese babel
\usepackage[english, vietnamese]{babel}
\usepackage{amssymb}
\usepackage{textcmds}
\usepackage{tabularray}
\usepackage{tabularx}
\usepackage{graphicx}
\usepackage{subcaption}
\usepackage{rotating} % để xoay bảng ngang
\usepackage{amssymb}  % để dùng \checkmark
\usepackage{pifont}   % để dùng dấu tick (✓) và dấu x

\usepackage{titletoc}
\renewcommand{\chaptername}{Chương}
\titlecontents{chapter}[0pt]{\addvspace{3pt}\bfseries}{\chaptername~\thecontentslabel\quad}{}{\hfill\contentspage}[\addvspace{5pt}]
\titlecontents{section}[2em]{\addvspace{2pt}}{\thecontentslabel\quad}{}{\titlerule*[0.5pc]{.}\contentspage}
\titlecontents{subsection}[4em]{\addvspace{0pt}}{\thecontentslabel\quad}{}{\titlerule*[0.5pc]{.}\contentspage}

\newcommand{\tick}{\ding{51}} % ✓
\newcommand{\ctick}{\raisebox{0pt}[1.8em][1.8em]{\centering\tick}}
\newcommand{\cross}{\ding{55}} % ✗ (nếu bạn cần dấu X)

%% ----- custom style----
\usepackage{amsmath}
\usepackage{nccmath}
\usepackage{tabularx}

\usepackage[numbers]{natbib}
\usepackage[noabbrev, nameinlink]{cleveref}
\usepackage{pdfpages}
\usepackage{enumitem}
\usepackage{textcomp}
\usepackage{pifont}
\usepackage{tabto}
\usepackage{booktabs}
\usepackage{array}
\usepackage{pdflscape}
\usepackage{forest}

\usepackage{xcolor}
\definecolor{codegreen}{rgb}{0,0.6,0}
\definecolor{codegray}{rgb}{0.5,0.5,0.5}
\definecolor{codepurple}{rgb}{0.58,0,0.82}
\definecolor{backcolour}{rgb}{0.95,0.95,0.92}

%% -- (start) image --
\usepackage{multirow}
\usepackage{graphicx}
\graphicspath{{Images/}}
%% -- (end)   image --


\usepackage{float}
\usepackage{wrapfig}
\usepackage{indentfirst}

%paragraph indentation
\setlength{\parindent}{2.5em} 
%paragraph spacing
\setlength{\parskip}{.5em}

\captionsetup[figure]{labelfont={small,bf},textfont={small,it},belowskip=-1pt,aboveskip=7pt}
% space remove between caption, figure, and text
\captionsetup[table]{labelfont={small,bf},textfont={small,it},belowskip=-1pt,aboveskip=7pt}
% space remove between caption, table, and text

\makeatletter
\renewcommand*\l@figure{\@dottedtocline{1}{1em}{3.2em}}
\makeatother

%%% define arrow
\newcommand{\myarrow}[1][1cm]{\mathrel{%
		\vcenter{\hbox{\rule[-1.5\fontdimen8\textfont3]{#1}{\fontdimen8\textfont3}}}%
		\mkern-4mu\hbox{\usefont{U}{lasy}{m}{n}\symbol{41}}}}

\newcommand{\localtextbulletone}{\textcolor{black}{\raisebox{.45ex}{\rule{.7ex}{.7ex}}}}
\renewcommand{\labelitemi}{\localtextbulletone}




%% ----- định lý, thuật toán (tên tiếng Việt) -----
\usepackage{longtable}
\usepackage{amsthm}
\usepackage{algorithm}
\usepackage{algpseudocode}
\floatname{algorithm}{Thuật toán}
\renewcommand{\listalgorithmname}{Danh sách thuật toán}
\algrenewcommand\algorithmicif{\textbf{nếu}}
\algrenewcommand\algorithmicthen{\textbf{thì}}
\algrenewcommand\algorithmicelse{\textbf{ngược lại}}
\algrenewcommand\algorithmicend{\textbf{hết}}
\algrenewcommand\algorithmicrequire{\textbf{Đầu vào:}}
\theoremstyle{plain}
\newtheorem{proposition}{Mệnh đề}[chapter]
\newtheorem{lemma}{Bổ đề}[chapter]
\theoremstyle{definition}
\newtheorem{definition}{Định nghĩa}[chapter]
\theoremstyle{remark}
\newtheorem{remark}{Nhận xét}[chapter]
\renewcommand{\proofname}{Chứng minh}

% Thêm macro riêng của báo cáo ở đây, ví dụ:
% \newcommand{\R}{\mathbb{R}}

\setlength{\textwidth}{160mm} % = 210mm - 30mm - 20mm
\setlength{\oddsidemargin}{4.6mm} % 30mm - 1in (from hoffset)
\setlength{\evensidemargin}{-5.4mm} % = 20mm - 1in (from hoffset)
\setlength{\topmargin}{-5.4mm} % = 0mm -1in + 20mm 
\setlength{\textheight}{231.9mm} % = 297mm -29.5mm -31.6mm - 14mm (12 to accomodate footline with pagenumber)
\setlength{\headheight}{16pt}

\usepackage{lmodern}
\usepackage{setspace} % increase interline spacing slightly
\setstretch{1.1}

\makeatletter
\setlength{\@fptop}{0pt}  % for aligning all floating figures/tables etc... to the top margin
\makeatother

\usepackage{booktabs}
\usepackage{microtype}
\usepackage{url}
\usepackage{mathtools}
\newcommand*\mystrut[1]{\vrule width0pt height0pt depth#1\relax}


\usepackage{xcolor}

\definecolor{codegreen}{rgb}{0,0.6,0}
\definecolor{codegray}{rgb}{0.5,0.5,0.5}
\definecolor{codepurple}{rgb}{0.58,0,0.82}
\definecolor{backcolour}{rgb}{0.95,0.95,0.92}




%%%%%% start - fancy %%%%%%%%%%%%%%%%%%%%%%%%
\usepackage{fancyhdr}
\renewcommand{\sectionmark}[1]{\markright{\thesection\ #1}}
\pagestyle{fancy}
	\fancyhf{}
	\renewcommand{\headrulewidth}{0pt}
	\renewcommand{\footrulewidth}{0pt}
	%% \fancyhead[OR]{\bfseries \nouppercase{\rightmark}} % current section: odd page, right side
	%% \fancyfoot[R]{\thepage}
	\fancyhead[R]{\bfseries \nouppercase{\rightmark}}
	% \fancyfoot[EL,OR]{\thepage} % page number: even/odd page, left/ right
    \fancyhead[C]{\thepage} % center + start page
\fancypagestyle{plain}{
	\fancyhf{}
	\renewcommand{\headrulewidth}{0pt}
	\renewcommand{\footrulewidth}{0pt}
	% \fancyfoot[EL,OR]{\thepage}}
	\fancyhead[C]{\thepage}} % center + start page

%%%%%% end - fancy %%%%%%%%%%%%%%%%%%%%%%%%%%
\usepackage{listings}
\lstset{language=[LaTeX]Tex,tabsize=4, basicstyle=\scriptsize\ttfamily, showstringspaces=false, numbers=left, numberstyle=\tiny, numbersep=10pt, breaklines=true, breakautoindent=true, breakindent=6pt}

\hypersetup{
    pdfborder={0 0 0},
	colorlinks=true,
	linkcolor=black,
	citecolor=blue,
	urlcolor=brown}
\urlstyle{same}

\makeatletter
\def\cleardoublepage{\clearpage\if@twoside \ifodd\c@page\else
    \hbox{}
    \thispagestyle{empty}
    \newpage
    \if@twocolumn\hbox{}\newpage\fi\fi\fi}
    
\makeatother \clearpage{\pagestyle{plain}\cleardoublepage}

%%%%% CHAPTER HEADER %%%%
\usepackage{color}
\usepackage{tikz}
\usetikzlibrary{calc,arrows.meta,decorations.pathreplacing,positioning}
\usepackage[explicit]{titlesec}
\newcommand*\chapterlabel{}
\newcommand{\chapterword}{Chương}

%\renewcommand{\thechapter}{\Roman{chapter}}
\titleformat{\chapter}[display]  % vị trí như mẫu gốc, không có hộp đen
	{\normalfont\bfseries\Huge} % format of the chapter
	{\xdef\chapterlabel{\chapterword~\thechapter}}     % the label with "Chương" prefix (xdef for immediate expansion)
 	{0pt} % separation between label and chapter-title
 	  {\begin{tikzpicture}[remember picture,overlay]
        \node[yshift=-8.2cm] at (current page.north west)
        {\begin{tikzpicture}[remember picture, overlay]
            \node[anchor=north east,yshift=-7.4cm,xshift=50mm,minimum height=30mm,inner sep=0mm] at (current page.north west)
            {\parbox[top][30mm][t]{25mm}{\raggedleft $\phantom{\textrm{l}}$\color{black}\mbox{\makebox[35.4mm][r]{\chapterlabel}}}};  % nhãn canh phải tại cùng một vị trí cho "Chương N" và "Phụ lục X"
            \node[anchor=north west,yshift=-7.2cm,xshift=67mm,text width=\textwidth,minimum height=30mm,inner sep=0mm] at (current page.north west)
                  {\parbox[top][30mm][t]{\textwidth}{\color{black}#1}};
           \end{tikzpicture}
          };
        \end{tikzpicture}
       \gdef\chapterlabel{}
      } % code before the title body

\newcounter{myparts}
\newcommand*\partlabel{}
\titleformat{\part}[display]  % type (section,chapter,etc...) to vary,  shape (eg display-type)
	{\normalfont\bfseries\Huge} % format of the part
	{\gdef\partlabel{\thepart\ }}     % the label 
 	{0pt} % separation between label and part-title
 	  {\setlength{\unitlength}{20mm}
	  \addtocounter{myparts}{1}
	  \begin{tikzpicture}[remember picture,overlay]
    \node[anchor=north west,xshift=-65mm,yshift=-6.9cm-\value{myparts}*20mm] at (current page.north east) % for unknown reasons: 3mm missing -> 65 instead of 62
      {\begin{tikzpicture}[remember picture, overlay]
        \draw[fill=black] (0,0) rectangle(62mm,20mm);   % -\value{myparts}\unitlength
        \node[anchor=north west,yshift=-6.1cm-\value{myparts}*20mm,xshift=-60.5mm,minimum height=30mm,inner sep=0mm] at (current page.north east)
        {\parbox[top][30mm][t]{55mm}{\raggedright \color{white}Part \partlabel $\phantom{\textrm{l}}$}};  %the phantom l is just to get better base-line alingement
        \node[anchor=north east,yshift=-6.1cm-\value{myparts}*20mm,xshift=-63.5mm,text width=\textwidth,minimum height=30mm,inner sep=0mm] at (current page.north east)
              {\parbox[top][30mm][t]{\textwidth}{\raggedleft \color{black}#1}};
       \end{tikzpicture}
      };
   \end{tikzpicture}
   \gdef\partlabel{}
  } % code before the title body
 
\titlespacing*{\chapter}{0pt}{35pt}{30pt}
\titlespacing*{\section}{0pt}{12pt}{*0}
\titlespacing*{\subsection}{0pt}{12pt}{*0}
\titlespacing*{\subsubsection}{0pt}{12pt}{*0}

%% --------------- define abbreviation ---------------
\newlist{abbrv}{itemize}{1}
\setlist[abbrv,1]{label=,labelwidth=1in,align=parleft,itemsep=0.1\baselineskip,leftmargin=!}
%% -----------()end) define abbreviation ---------------

% % Headnote
% \fancyhead{} % clear all header fields
% \fancyhead[L]{
% 	\begin{tabular}{rl}
% 		\begin{picture}(25,15)(0,0)
% 			\put(0,-8){\includegraphics[width=8mm, height=8mm]{Images/hcmut.png}}
% 			%\put(0,-8){\epsfig{width=10mm,figure=hcmut.eps}}
% 		\end{picture}&
% 		%\includegraphics[width=8mm, height=8mm]{hcmut.png} & %
% 		\begin{tabular}{l}
% 			\textbf{\bf \ttfamily Trường Đại Học Bách Khoa Tp.Hồ Chí Minh}\\
% 			\textbf{\bf \ttfamily Khoa Khoa Học \& Kỹ Thuật Máy Tính}
% 		\end{tabular} 	
% 	\end{tabular}
% }
% \fancyhead[R]{
% 	\begin{tabular}{l}
% 		\tiny \bf \\
% 		\tiny \bf 
% \end{tabular}  }

% % Footnote
% \fancyfoot{} % clear all footer fields
% \fancyfoot[L]{\scriptsize \ttfamily  Đồ án chuyên ngành Khoa Học Máy Tính}
% \fancyfoot[R]{\scriptsize \ttfamily Trang {\thepage}/\pageref{LastPage}}
% \renewcommand{\headrulewidth}{0.3pt}
% \renewcommand{\footrulewidth}{0.3pt}

%%%%%%%%%%%%%%% define JSON format %%%%%%%%%%%%%%%
\colorlet{punct}{red!60!black}
\definecolor{background}{HTML}{EEEEEE}
\definecolor{delim}{RGB}{20,105,176}
\colorlet{numb}{magenta!60!black}

\lstdefinelanguage{json}{
    basicstyle=\ttfamily\footnotesize,
    numbers=left,
    numberstyle=\scriptsize,
    stepnumber=1,
    numbersep=8pt,
    showstringspaces=false,
    breaklines=true,
    keepspaces=true,
    breakatwhitespace=false,
    showtabs=false
}

%%%%%%%%%%%%%%% end define JSON format %%%%%%%%%%%%%%%

\usepackage{lastpage}

%%%%%%%%%%%%% (start) add package %%%%%%%%%%%%%%%
\usepackage{graphicx}
\usepackage{cases}
\usepackage{makecell}
\usepackage{times}
%%%%%%%%%%%%% (end) add package %%%%%%%%%%%%%%%%%
\begin{document}

\lstdefinestyle{mystyle}{
    backgroundcolor=\color{backcolour},   
    commentstyle=\color{codegreen},
    keywordstyle=\color{magenta},
    numberstyle=\tiny\color{codegray},
    stringstyle=\color{codepurple},
    basicstyle=\ttfamily\footnotesize,
    breakatwhitespace=false,         
    breaklines=true,                 
    captionpos=b,                    
    keepspaces=true,                 
    numbers=left,                    
    numbersep=5pt,                  
    showspaces=false,                
    showstringspaces=false,
    showtabs=false,                  
    tabsize=2
}

\lstset{style=mystyle}


\onehalfspacing
%%%%%%%%%%%%%%%%%%%%%%%%%%%%%%%%%%%%%%%%%%%%%%
%%%%% HEAD: Book-Begin
%%%%%%%%%%%%%%%%%%%%%%%%%%%%%%%%%%%%%%%%%%%%%%
\frontmatter
    \begin{titlepage}
\thispagestyle{empty}

% --- KHUNG VIỀN ---
\begin{center}
    \begin{tikzpicture}[remember picture, overlay]
        \draw[line width = 3pt] ($(current page.north west) + (3.0cm,-2.0cm)$) rectangle ($(current page.south east) + (-2.0cm,2.0cm) $);
        \draw[line width = 1pt] ($(current page.north west) + (2.9cm,-1.9cm)$) rectangle ($(current page.south east) + (-1.9cm,1.9cm) $);
    \end{tikzpicture}

% --- PHẦN HEADER ---
    \begin{large}
        \textbf{ĐẠI HỌC QUỐC GIA TP.HCM} \\
        \textbf{TRƯỜNG ĐẠI HỌC BÁCH KHOA} \\
        \textbf{KHOA KHOA HỌC \& KỸ THUẬT MÁY TÍNH}
    \end{large} \\

    \vspace{0.5cm}

    \includegraphics[scale=0.3]{Images/hcmut.png}\\

    \vspace{0.5cm}

    % BÁO CÁO THỰC TẬP / ĐỀ CƯƠNG LUẬN VĂN / LUẬN VĂN THẠC SĨ / ĐỒ ÁN ...
    {\fontsize{16pt}{1}\selectfont \textbf{[LOẠI BÁO CÁO]}}\\[0.75cm]

% --- TÊN ĐỀ TÀI ---
    \begin{minipage}{0.85\textwidth}
        \centering\textbf{\large [TÊN ĐỀ TÀI]}
    \end{minipage}\\[0.4cm]
    {\large\textit{[Phụ đề, có thể bỏ]}}\\[0.75cm]

    {\fontsize{16pt}{1}\selectfont \textbf{NGÀNH: [TÊN NGÀNH]}}

    \vfill

% --- BẢNG THÔNG TIN ---
    \renewcommand{\arraystretch}{1.3}
    \begin{tabular}{r c l}
        \textbf{\large Giảng viên hướng dẫn} & : & \large [Học hàm, học vị, họ tên] \\
                                           &   & \large [Đồng hướng dẫn, nếu có] \\ [0.3cm]
        \textbf{\large Học viên thực hiện}   & : & \large [Họ tên] -- [MSHV] \\
    \end{tabular}

    \vfill

% --- FOOTER ---
    {\fontsize{14pt}{1}\selectfont TP. HỒ CHÍ MINH, THÁNG [MM/YYYY]}

\end{center}
\end{titlepage}

    
    \pagenumbering{roman}
    \setcounter{page}{1}
    
    \chapter*{Lời cảm ơn}
\markboth{Lời cảm ơn}{Lời cảm ơn}
\addcontentsline{toc}{chapter}{Lời cảm ơn}

\onehalfspacing
\vspace{1.0cm}

[Lời cảm ơn gửi giảng viên hướng dẫn, Khoa, gia đình và bạn bè.] \par

\begin{flushright}
\begin{tabular}{@{}c@{}}
    TP. Hồ Chí Minh, tháng [MM] năm [YYYY]\\
    \textbf{Học viên thực hiện}\\ [2cm]
    \textbf{[Họ tên]}
\end{tabular}
\end{flushright}

\bigskip

    \chapter*{Tóm tắt}
\markboth{Tóm tắt}{Tóm tắt}
\addcontentsline{toc}{chapter}{Tóm tắt}
\onehalfspacing
\vspace{1.0cm}

[Bối cảnh và vấn đề.] [Phương pháp đề xuất.] [Cách đánh giá.] [Kết quả chính.]

\bigskip

    \chapter*{Abstract}
\markboth{Abstract}{Abstract}
\addcontentsline{toc}{chapter}{Abstract}
\onehalfspacing
\vspace{1.0cm}

\begin{otherlanguage}{english}
[Context and problem.] [Proposed method.] [Evaluation.] [Main results.]
\end{otherlanguage}

\bigskip

    \chapter*{Lời cam đoan}
\markboth{Lời cam đoan}{Lời cam đoan}
\addcontentsline{toc}{chapter}{Lời cam đoan}
\onehalfspacing
\vspace{1.0cm}

Tôi xin cam đoan nội dung trong báo cáo này là kết quả tôi thực hiện dưới sự hướng dẫn của [giảng viên hướng dẫn], ngoại trừ các nội dung đã được trích dẫn rõ nguồn tham khảo. [Bổ sung cam kết về số liệu và công cụ hỗ trợ, nếu cần.] \par

\begin{flushright}
\begin{tabular}{@{}c@{}}
    TP. Hồ Chí Minh, tháng [MM] năm [YYYY]\\
    \textbf{Học viên thực hiện}\\ [2cm]
    \textbf{[Họ tên]}
\end{tabular}
\end{flushright}

\bigskip

    
    \pagestyle{plain}

    % Đưa vào TOC trước khi in TOC để lấy đúng trang bắt đầu
    \cleardoublepage
    \phantomsection
    \addcontentsline{toc}{chapter}{Mục lục} 
    \tableofcontents
    \cleardoublepage
 
    \phantomsection
    \addcontentsline{toc}{chapter}{Danh sách hình vẽ} 
    \listoffigures
    \cleardoublepage
 
    \phantomsection
    \addcontentsline{toc}{chapter}{Danh sách bảng} 
    \listoftables
    \cleardoublepage
    
    \phantomsection
    \addcontentsline{toc}{chapter}{Thuật ngữ \& từ viết tắt} 
    \chapter*{Thuật ngữ \& từ viết tắt}
\vspace{1.0cm}
\begin{abbrv}
    \item[AI] Artificial Intelligence -- trí tuệ nhân tạo
    \item[LLM] Large Language Model -- mô hình ngôn ngữ lớn
    \item[{[VT]}] [Tên đầy đủ] -- [nghĩa tiếng Việt]
\end{abbrv}

    \if@openright\cleardoublepage\else\clearpage\fi


\setlength{\parskip}{1em}

%%%%%%%%%%%%%%%%%%%%%%%%%%%%%%%%%%%%%%%%%%%%%%
%%%%% MAIN: The chapters and appendix of the thesis
%%%%%%%%%%%%%%%%%%%%%%%%%%%%%%%%%%%%%%%%%%%%%%
\mainmatter
    \pagestyle{fancy}
    \chapter{Giới thiệu}\label{chapter:introduction}
\pagestyle{fancy}

\section{Bối cảnh và động lực}
[Nội dung.] Trích dẫn mẫu: \cite{sample2024}.

\section{Mục tiêu}
\begin{itemize}
    \item Mục tiêu 1: [nội dung].
    \item Mục tiêu 2: [nội dung].
\end{itemize}

\section{Phạm vi và giới hạn}
[Nội dung.]

\section{Cấu trúc báo cáo}
Chương~\ref{chapter:background} trình bày cơ sở lý thuyết, Chương~\ref{chapter:method} mô tả phương pháp, Chương~\ref{chapter:experiment} trình bày thực nghiệm và Chương~\ref{chapter:conclusion} kết luận.
 % Giới thiệu
        \newpage
    \chapter{Cơ sở lý thuyết và công trình liên quan}\label{chapter:background}
\pagestyle{fancy}

\section{Kiến thức nền tảng}

\begin{definition}[Tên định nghĩa]\label{def:sample}
[Phát biểu định nghĩa.]
\end{definition}

\begin{proposition}\label{prop:sample}
Với mọi $x \in \mathbb{R}$, $x^2 \ge 0$.
\end{proposition}

\begin{proof}
[Chứng minh.]
\end{proof}

\section{Công trình liên quan}
Bảng~\ref{tab:related} là mẫu bảng so sánh.

\begin{table}[ht]
  \centering
  \caption{Mẫu bảng so sánh.}\label{tab:related}
  \begin{tabular}{@{}lccc@{}}
    \toprule
    Công trình & Tiêu chí A & Tiêu chí B & Tiêu chí C \\
    \midrule
    {[Công trình 1]} & \tick & \cross & \tick \\
    {[Công trình 2]} & \cross & \tick & \tick \\
    Đề tài này     & \tick & \tick & \tick \\
    \bottomrule
  \end{tabular}
\end{table}
 % Cơ sở lý thuyết và công trình liên quan
        \newpage
    \chapter{Phương pháp đề xuất}\label{chapter:method}
\pagestyle{fancy}

\section{Tổng quan}
Hình~\ref{fig:sample} là mẫu hình.

\begin{figure}[ht]
  \centering
  \includegraphics[width=0.25\textwidth]{hcmut.png}
  \caption{Mẫu hình (thay bằng hình của đề tài, đặt trong thư mục \texttt{Images/}).}\label{fig:sample}
\end{figure}

\section{Thuật toán}
Thuật toán~\ref{alg:sample} là mẫu thuật toán.

\begin{algorithm}[ht]
  \caption{Mẫu thuật toán}\label{alg:sample}
  \begin{algorithmic}[1]
    \Require Dãy số $a_1, \dots, a_n$
    \State $s \gets 0$
    \For{$i \gets 1$ \textbf{to} $n$}
      \If{$a_i > 0$}
        \State $s \gets s + a_i$
      \EndIf
    \EndFor
    \State \Return $s$
  \end{algorithmic}
\end{algorithm}
 % Phương pháp đề xuất
        \newpage
    \chapter{Thực nghiệm và kết quả}\label{chapter:experiment}
\pagestyle{fancy}

\section{Thiết lập thực nghiệm}
[Dữ liệu, độ đo, các phương pháp so sánh.]

\section{Kết quả}
[Kết quả chính.]

\begin{lstlisting}[language=Python, caption={Mẫu mã nguồn.}]
def add(a, b):
    return a + b
\end{lstlisting}

\section{Thảo luận}
[Nội dung.]
 % Thực nghiệm và kết quả
        \newpage
    \chapter{Kết luận}\label{chapter:conclusion}
\pagestyle{fancy}

\section{Kết quả đạt được}
[Nội dung.]

\section{Hạn chế và hướng phát triển}
[Nội dung.]
 % Kết luận
        \newpage


    %% ----- Phụ lục -----
    \appendix
    \renewcommand{\chapterword}{Phụ lục}
    \addtocontents{toc}{\protect\renewcommand{\protect\chaptername}{Phụ lục}}
\chapter{Bảng ký hiệu}\label{appendix:notation}
\pagestyle{fancy}

\begin{longtable}{@{}p{0.20\textwidth} p{0.74\textwidth}@{}}
\caption{Các ký hiệu dùng trong báo cáo.}\label{tab:notation}\\
\toprule
Ký hiệu & Ý nghĩa \\
\midrule
\endfirsthead
\toprule
Ký hiệu & Ý nghĩa \\
\midrule
\endhead
\bottomrule
\endfoot
$n$ & [Ý nghĩa] \\
$\mathbb{R}$ & Tập số thực \\
\end{longtable}

        \newpage

\backmatter
    \bibliographystyle{references/IEEEtran}
    \bibliography{references/refs}
    \addcontentsline{toc}{chapter}{Tài liệu tham khảo}
    % \chapter*{Phần lý lịch trích ngang}
\addcontentsline{toc}{chapter}{Phần lý lịch trích ngang}
\pagestyle{fancy}

% LƯU Ý: các mục trong ngoặc vuông là chỗ giữ chỗ, điền trước khi nộp.
\vspace{0.5cm}
\begin{itemize}
    \item[] Họ và tên: [Họ tên]
    \vspace{0.2cm}
    \item[] Mã số học viên: [MSHV]
    \vspace{0.2cm}
    \item[] \makebox[0.6\linewidth][l]{Ngày, tháng, năm sinh: [điền]} Nơi sinh: [điền]
    \vspace{0.2cm}
    \item[] Địa chỉ liên lạc: [điền]
\end{itemize}

\section*{Quá trình đào tạo}
\begin{itemize}
    \item \textbf{[thời gian]}: [bậc đại học, ngành, trường].

    \item \textbf{[thời gian] -- nay}: Học viên cao học ngành [tên ngành], Khoa Khoa học \& Kỹ thuật Máy tính, Trường Đại học Bách Khoa -- ĐHQG TP.HCM.
\end{itemize}

\section*{Quá trình công tác}
\begin{itemize}
    \item \textbf{[thời gian]}: [vị trí, đơn vị công tác].
\end{itemize}
 % lý lịch trích ngang (luận văn)


\end{document}

